\paragraph{Multimodal sentiment fusion.}
Early work on \textsc{CMU-MOSI} and \textsc{CMU-MOSEI} aligned
modality-specific recurrent encoders through tensor / outer-product
fusion (TFN \citep{zadeh2017tfn}, LMF \citep{liu2018lmf}, MFN
\citep{zadeh2018mfn}, Graph-MFN \citep{zadeh2018gmfn}).
The transformer era replaced these with cross-modal attention
\citep{tsai2019mult} and BERT-grounded fusion (MAG-BERT
\citep{rahman2020magbert}, MISA \citep{hazarika2020misa}, Self-MM
\citep{yu2021selfmm}, MMIM \citep{han2021mmim}).
A second line of work disentangles modality-shared and
modality-specific subspaces (MFM \citep{tsai2019mfm}, DMD
\citep{li2023dmd}, EMOE \citep{emoe2024}); a third focuses on
cross-modal alignment through contrastive or denoising objectives
\citep{wu2023dbf,boost2024}.

\paragraph{Information bottleneck for multimodal learning.}
The IB principle \citep{tishby2000information,tishby2015dnn} compresses
input $X$ into a bottleneck $B$ by minimising $I(X;B)$ while maximising
$I(B;Y)$.
Its variational form (VIB) \citep{alemi2017vib,kingma2014vae} makes the
objective tractable for deep networks.
ITHP \citep{ithp2024} introduces a hierarchical IB chain for
multimodal perception, while CyIN cyclically translates between
modality bottlenecks; CaMIB \citep{camib2026} adds causal-aware
ingredients to the same recipe.
C-MIB \citep{cmib2023} and DLF \citep{dlf2024} regularise cross-modal
mutual information.
\methodname{} differs from this family in two ways.
First, while existing IB methods use $\log\sigma^{2}$ only inside the
KL term of the VIB objective, we expose it as a \emph{first-class
modulation signal} for downstream attention and gating.
Second, IB methods that perform translation
(CyIN, ITHP) require auxiliary cross-modal decoders even at inference;
\methodname{} does not perform any cross-modal translation and
therefore has a strictly smaller inference-time graph.

\paragraph{Routing and dynamic primary selection.}
MODS \citep{wang2023cross} is the closest predecessor.
It introduces a Modality Selector that uses Gumbel-softmax to pick a
per-sample primary modality, and an equal-weight cross-attention
fusion (PCCA).
\methodname{} reuses the MSelector backbone, but adds the IB-guided
modulation of \S\ref{sec:method:ibca}--\ref{sec:method:gate}, and
supervises the selector in stage 2 with a routing-IB loss that uses
the per-modality task error rather than a fixed text prior.
DEVA \citep{deva2024}, KuDA \citep{kuda2024}, MOAC \citep{moac2025}
and TMSON \citep{tmson2024} are concurrent dynamic-routing or
adaptive-fusion methods compared in Table~\ref{tab:main_mosi_mosei};
none of them couples routing to a per-token IB confidence.
Adaptive Information Gates (Eq.~\ref{eq:adagate}) are also reminiscent
of the gated fusion in MAG-BERT and DBF \citep{wu2023dbf}, but
condition on the IB-compressed bottleneck rather than on raw
features.
